% Part: proof-theory
% Chapter: sequent-calculus
% Section: proof-examples

\documentclass[../../../include/open-logic-section]{subfiles}

\begin{document}

\olfileid{pt}{seq}{ex}

\olsection{Examples of Proofs}

To give !!{proof}s of sequents it is typically convenient to work your
way upwards from the end-sequent by selecting !!a{formula} in it as
the active !!{formula} and writing the premises of the corresponding
rule on top of the sequent, then repeating the process until it leads
to axioms. For instance, suppose we want to prove the sequent
\[(!C \land !D) \lif !E \Sequent !C \lif (!D \lif !E)\]
Consider $!C \lif (!D \lif !E)$: it is of the form $!A \lif !B$ and
occurs on the right of the sequent. Matching the entire sequent with
the $\RightR{\lif}$ rule of \Log{G1c},
\[\Axiom$ !A, \Gamma \fCenter \Delta, !B$
\RightLabel{\RightR{\lif}}
\UnaryInf$ \Gamma \fCenter \Delta, !A \lif !B $
\DisplayProof\]
suggests we should take $\Gamma = \{!C \lif (!D \lif !E)\}$, $\Delta =
\emptyset$, $!A \ident !C$ and $!B \ident !D \lif !E$. So the last
inference of a proof could be
\[\Axiom$ !C, (!C \land !D) \lif !E \fCenter !D \lif !E$
\RightLabel{\RightR{\lif}}
\UnaryInf$(!C \land !D) \lif !E \fCenter !C \lif (!D \lif !E)$
\DisplayProof\]
If we repeat this idea, focusing on $!D \lif !E$ as active formula, we
generate
\[\Axiom$ !D, !C, (!C \land !D) \lif !E \fCenter !E$
\RightLabel{\RightR{\lif}}
\UnaryInf$!C, (!C \land !D) \lif !E \fCenter !D \lif !E$
\RightLabel{\RightR{\lif}}
\UnaryInf$(!C \land !D) \lif !E \fCenter !C \lif (!D \lif !E)$
\DisplayProof\]
Now we have to focus on $(!C \land !D) \lif !E$ and use the $\LeftR{\lif}$ rule of \Log{G1c},
\[\Axiom$ \Gamma \fCenter \Delta, !A$
\Axiom$ !B, \Gamma \fCenter \Delta$
\RightLabel{\LeftR{\lif}}
\BinaryInf$ !A \lif !B, \Gamma \fCenter \Delta$
\DisplayProof\]
This leads us to
\[\Axiom$!D, !C \fCenter !E, !C \land !D$
\Axiom$!D, !C, !E \fCenter !E$
\RightLabel{\LeftR{\lif}}
\BinaryInf$ !D, !C, (!C \land !D) \lif !E \fCenter !E$
\RightLabel{\RightR{\lif}}
\UnaryInf$!C, (!C \land !D) \lif !E \fCenter !D \lif !E$
\RightLabel{\RightR{\lif}}
\UnaryInf$(!C \land !D) \lif !E \fCenter !C \lif (!D \lif !E)$
\DisplayProof\]
Using the !!{formula} $!C \land !D$ in the left topmost sequent as
principal and the $\RightR{\land}$ rule gives us
\[
\Axiom$!D, !C \fCenter !E, !C$
\Axiom$!D, !C \fCenter !E, !D$
\RightLabel{\RightR{\land}}
\BinaryInf$!D, !C \fCenter !E, !C \land !D$
\Axiom$!D, !C, !E \fCenter !E$
\RightLabel{\RightR{\lif}}
\BinaryInf$ !D, !C, (!C \land !D) \lif !E \fCenter !E$
\RightLabel{\RightR{\lif}}
\UnaryInf$!C, (!C \land !D) \lif !E \fCenter !D \lif !E$
\RightLabel{\RightR{\lif}}
\UnaryInf$(!C \land !D) \lif !E \fCenter !C \lif (!D \lif !E)$
\DisplayProof
\]
So far so good. Note that everything we did also works in \Log{G3c},
since the rules we've made use of so far are the same in \Log{G1c}
and~\Log{G3c}. We've arrived at !!a{proof} fragment where all top
sequents contain the same !!{formula} on the left as on the right. But
they are not quite axioms yet.

In \Log{G1c}, we have to give !!{proof}s of the topmost sequents from
axioms of the form $!A \Sequent !A$. This is easily done using the
weakening rules, e.g., the leftmost topmost sequent $!D, !C \Sequent
!E, !C$ can be !!{prove}d
as follows:
\[
\Axiom$!C \fCenter !C$
\RightLabel{\LeftR{\Weakening}}
\UnaryInf$!D, !C \fCenter !C$
\RightLabel{\RightR{\Weakening}}
\UnaryInf$!D, !C \fCenter !E, !C$
\]
Overall, we have !!a{proof} in \Log{G1c} which looks like this:
\[
\Axiom$!C \fCenter !C$
\RightLabel{\LeftR{\Weakening}}
\UnaryInf$!D, !C \fCenter !C$
\RightLabel{\RightR{\Weakening}}
\UnaryInf$!D, !C \fCenter !E, !C$
\Axiom$!D \fCenter !D$
\RightLabel{\LeftR{\Weakening}}
\UnaryInf$!D, !C \fCenter !D$
\RightLabel{\RightR{\Weakening}}
\UnaryInf$!D, !C \fCenter !E, !D$
\RightLabel{\RightR{\land}}
\BinaryInf$!D, !C \fCenter !E, !C \land !D$
\Axiom$!E \fCenter !E$
\RightLabel{\LeftR{\Weakening}}
\UnaryInf$!C, !E \fCenter !C$
\RightLabel{\LeftR{\Weakening}}
\UnaryInf$!D, !C, !E \fCenter !E$
\RightLabel{\LeftR{\lif}}
\BinaryInf$ !D, !C, (!C \land !D) \lif !E \fCenter !E$
\RightLabel{\RightR{\lif}}
\UnaryInf$!C, (!C \land !D) \lif !E \fCenter !D \lif !E$
\RightLabel{\RightR{\lif}}
\UnaryInf$(!C \land !D) \lif !E \fCenter !C \lif (!D \lif !E)$
\DisplayProof
\]
The structure of this proof (in particular, its height) is independent
of what the !!{formula}s $!C$, $!D$, and~$!E$ are. We can replace
these symbols in the !!{proof} above by any !!{formula}s and the
result will still be a !!{proof} in~\Log{G1c}. The proof in \Log{G1c}
is \emph{schematic}.

In \Log{G3c}, axioms are sequents of the form $!A, \Gamma \Sequent
\Delta, !A$ where $!A$ must be atomic. So although $!D, !C \Sequent
!E, !C$ looks like an axiom of \Log{G3c} ($\Gamma = \{!D\}$; $\Delta =
\{!E\}$), it is only \emph{really} an axiom if $!C$ in fact is atomic.
Our proof fragment is only a complete proof in \Log{G3c} if $!C$,
$!D$, and $!E$ are atomic !!{formula}s.

Although we haven't, strictly speaking, shown that
\[\Log{G1c} \Proves (!C \land !D) \lif !E \fCenter !C \lif (!D \lif
!E),\] in practice this is all we need to do (and can do). That's
because every sequent of the form $!A, \Gamma \Sequent \Delta, !A$ is
provable in \Log{G3c}, even if $!A$ itself is not atomic. We can show
this by induction on the depth~$\depth{!A}$ of~$!A$.

\begin{defn}\ollabel{defn:depth}
The \emph{depth}~$\depth{!A}$ of !!a{formula} is defined as follows:
\begin{enumerate}
  \item If $!A$ is atomic, $\depth{!A} = 0$.
  \item If $!A \ident \lnot !B$, $!A \ident \lforall[x][A]$, or 
  $!A \ident \lexists[x][A]$, then $\depth{!A} = \depth{!B} + 1$.
  \item If $!A \ident (!B \land !C)$, $(!B \lor !C)$, or $(!B \lif
  !C)$, then $\depth{!A} = \max(\depth{!B},\depth{!C}) + 1$.
\end{enumerate}
\end{defn}

It is not hard to prove that $\depth{A(x)} = \depth{A(t)}$ for any
term~$t$. The important fact for us is that in any rule, the depth of
the principal !!{formula} is always greater than the depth of the
active !!{formula}s. For instance, we have:
\begin{align*}
\depth{P(x)} = \depth{P(a)} & = 0,\\
\depth{\lforall[x][P(x)]} & = 1,\\
\depth{(\lforall[x][P(x)] \lif P(a))} & = \max(1,0) + 1 = 2.
\end{align*}
We can use this to prove some results by induction on~$\depth{!A}$,
such as the following.

\begin{prop}\ollabel{prop:G1c-axioms}
For any~$!A$, there is !!a{proof} in $\Log{G1c}$ of $!A \Sequent !A$
in which all axioms are atomic.
\end{prop}

\begin{proof}
  By induction on~$\depth{!A}$. If $\depth{!A} = 0$ then $!A$ is
  atomic and $!A \Sequent !A$ already is a !!{proof} in~\Log{G1c} of
  the required form. If $\depth{!A} > 0$, $!A$ is constructed from
  !!{formula}s of lower depth, e.g, $!A \ident \lnot !B$, $!A \ident
  (!B \land !C)$, or $!A \ident \lforall[x][A(x)]$. The inductive
  hypothesis is that for all $!D$ with $\depth{!D} < \depth{!A}$,
  $\Log{G1c} \Proves !D \Sequent !D$ from atomic axioms.

  If $!A \ident \lnot !B$, then $\depth{!B} < \depth{!A}$ and by
  inductive hypothesis, there is !!a{proof}~$\pi_1$ in \Log{G1c} from
  atomic axioms of
  $!B \Sequent !B$. We can extend this to !!a{proof}
  of $\lnot !B \Sequent \lnot !B$ as follows:
  \[
  \AxiomC{}
  \RightLabel{$\pi_1$}
  \Deduce$!B \fCenter !B$
  \RightLabel{\LeftR{\lnot}}
  \UnaryInf$\lnot !B, !B \fCenter $
  \RightLabel{\RightR{\lnot}}
  \UnaryInf$\lnot !B \fCenter \lnot !B$
  \DisplayProof
\]

  If $!A \ident (!B \land !C)$, then $\depth{!B} < \depth{!A}$ and
  $\depth{!C} < \depth{!A}$. By inductive hypothesis, there are
  !!{proof}s~$\pi_1$ and~$\pi_2$ in \Log{G1c} from atomic axioms of $!B
  \Sequent !B$ and $!C \Sequent !C$. We can extend these
  to !!a{proof} of $!B \land !C \Sequent !B \land !C$
  as follows:
  \[
  \AxiomC{}
  \RightLabel{$\pi_1$}
  \Deduce$!B, \fCenter !B$
  \RightLabel{\LeftR{\Weakening}}
  \UnaryInf$!B, !C, \fCenter !B$
  \RightLabel{\LeftR{\land}}
  \UnaryInf$!B \land !C, !C\fCenter !B$
  \RightLabel{\LeftR{\land}}
  \UnaryInf$!B \land !C, !B \land !C \fCenter !B$
  \AxiomC{}
  \RightLabel{$\pi_2$}
  \Deduce$!C, \fCenter !C$
  \RightLabel{\LeftR{\Weakening}}
  \UnaryInf$!B, !C, \fCenter !C$
  \RightLabel{\LeftR{\land}}
  \UnaryInf$!B \land !C, !C \fCenter !C$
  \RightLabel{\LeftR{\land}}
  \UnaryInf$!B \land !C, !B \land !C \fCenter !C$
  \RightLabel{\RightR{\land}}
  \BinaryInf$!B \land !C \fCenter !B \land !C$
  \RightLabel{\LeftR{\Contraction}}
  \UnaryInf$!B \land !C \fCenter !B$
  \DisplayProof
  \]

  Now suppose $!A \ident \lforall[x][B(x)]$. Take !!a{constant}~$c$
  that does not occur in $\lforall[x][B(x)]$. Then $\depth{B(c)} =
  \depth{!B(x)} < \depth{!A}$ and by inductive hypothesis, there is
  !!a{proof}~$\pi_1$ in \Log{G1c} from atomic axioms of $!B(c)
  \Sequent !B(c)$. We can extend this to !!a{proof} of
  $\lforall[x][B(x)] \Sequent \lforall[x][B(x)]$ as
  follows:
  \[
  \AxiomC{}
  \RightLabel{$\pi_1$}
  \Deduce$!B(c) \fCenter !B(c)$
  \RightLabel{\LeftR{\lforall}}
  \UnaryInf$\lforall[x][!B(x)] \fCenter !B(c)$
  \RightLabel{\RightR{\lforall}}
  \UnaryInf$\lforall[x][!B(x)] \fCenter \lforall[x][!B(x)]$
  \DisplayProof
\]

  The remaining cases are left as exercises.
\end{proof}

\begin{prob}
  Complete the proof of \olref{prop:G1c-axioms} by carrying out the
  cases where $!A \ident (!B \lor !C)$, $!A \ident (!B \lif !C)$ and
  $!A \ident \lexists[x][B(x)]$.
\end{prob}

For our purposes, we could also have used the !!{proof}
\[
\AxiomC{}
\RightLabel{$\pi_1$}
\Deduce$!B \fCenter !B$
\RightLabel{\RightR{\lnot}}
\UnaryInf$\fCenter !B, \lnot !B$
\RightLabel{\LeftR{\lnot}}
\UnaryInf$\lnot !B \fCenter \lnot !B$
\DisplayProof
\]
in the case where $!A \ident \lnot !B$. However, there are sequent
systems in which this wouldn't have worked. The intuitionistic
sequent calculus \Log{G1i} is like \Log{G1c}, except it restricts all
sequents to at most one !!{formula} on the right side of a sequent. If
$\Delta = \emptyset$, then the first !!{proof} would also be
!!a{proof} in~\Log{G1i}, but the second would not be, since
one sequent contains both $!B$ and~$\lnot !B$ on the right.

Note that even in \Log{G1c} we could not have used an alternate order
of rules in the case where $!A \ident \lforall[x][!B(x)]$. The
following is \emph{not} !!a{proof}
  \[
  \AxiomC{}
  \RightLabel{$\pi_1$}
  \Deduce$!B(c) \fCenter !B(c)$
  \RightLabel{\RightR{\lforall}}
  \UnaryInf$!B(c) \fCenter \lforall[x][!B(x)]$
  \RightLabel{\LeftR{\lforall}}
  \UnaryInf$\lforall[x][!B(x)] \fCenter \lforall[x][!B(x)]$
  \DisplayProof
\]
since the eigenvariable condition on~$\RightR{\lforall}$ is violated.

\begin{prop}\ollabel{prop:G3c-axioms}
Any sequent of the form $!A, \Gamma
\Sequent \Delta, !A$ (with $!A$ \emph{not} necessarily atomic) is
provable in~\Log{G3c}.
\end{prop}

\begin{proof}
The proof is very similar to the one of \olref{prop:G1c-axioms}, but
we must be careful with the inductive hypothesis. The case where $n =
\depth{!A} = 0$ is again easy. If $n=0$, $!A$ is atomic, and $!A,
\Gamma \Sequent \Delta, !A$ is an axiom of~\Log{G3c}.

Now suppose $n>0$. Again we distinguish cases. The inductive
hypothesis is: for all $D!$ with $\depth{!D} < n$, and all $\Pi$
and $\Lambda$, $\Log{G3c} \Proves !D, \Pi \Sequent \Lambda, !D$. Here
is the case for $!A \ident (!B \land !C)$. We use the inductive
hypothesis twice: once for $!D = !B$, $\Pi = !C, \Gamma$, and $\Lambda
= \Delta$ to get !!a{proof}~$\pi_1$ of $!B, !C, \Gamma \Sequent
\Delta, !B$, and once for $!D = !C$, $\Pi = !B, \Lambda$, and $\Lambda
= \Delta$ to get !!a{proof}~$\pi_2$ of $!B, !C, \Gamma \Sequent
\Delta, !C$. Here is the !!{proof} of~$!B \land !C, \Gamma \Sequent
\Delta, !B \land !C$:
\[
\AxiomC{}
\RightLabel{$\pi_1$}
\Deduce$!B, !C, \Gamma \fCenter \Delta, !B$
\RightLabel{\LeftR{\land}}
\UnaryInf$!B \land !C, \Gamma \fCenter \Delta, !B$
\AxiomC{}
\RightLabel{$\pi_2$}
\Deduce$!B, !C, \Gamma \fCenter \Delta, !C$
\RightLabel{\LeftR{\land}}
\UnaryInf$!B \land !C, \Gamma \fCenter \Delta, !C$
\RightLabel{\RightR{\land}}
\BinaryInf$!B \land !C, \Gamma \fCenter \Delta, !B \land !C$
\DisplayProof
\]

For $!A \ident \lforall[x][!B(x)]$, pick !!a{constant} not
in~$\lforall[x][!B(x)]$, $\Gamma$, or~$\Delta$. Apply the inductive
hypothesis for $!D = !B(c)$, $\Pi = \lforall[x][!B(x)], \Gamma$
and $\Lambda = \Delta$. This gives !!a{proof}~$\pi_1$ of $!B(c),
\lforall[x][!B(x)], \Gamma \Sequent \Delta, !B(c)$, from which we
obtain:
\[
\AxiomC{}
\RightLabel{$\pi_1$}
\Deduce$!B(c), \lforall[x][!B(x)], \Gamma \fCenter \Delta, !B(c)$
\RightLabel{\LeftR{\lforall}}
\UnaryInf$\lforall[x][!B(x)], \Gamma \fCenter \Delta, !B(c)$
\RightLabel{\RightR{\lforall}}
\UnaryInf$\lforall[x][!B(x)], \Gamma \fCenter \Delta, \lforall[x][!B(x)]$
\DisplayProof
\]
The eigenvariable condition on $\RightR{\lforall}$ is satisfied by the
way we chose~$c$.
\end{proof}

\begin{prob}
  Continue the proof of \olref{prop:G3c-axioms} by carrying out the
  cases where $!A \ident (!B \lif !C)$ and $!A \ident
  \lexists[x][B(x)]$.
\end{prob}

\end{document}